%% Supplementary material for main.tex. Compile main.tex first (its .aux
%% supplies the cross-referenced equation numbers), then:
%%   latexmk -pdf supplement.tex
\documentclass[10pt,journal]{IEEEtran}

\usepackage{amsmath,amssymb,amsfonts,amsthm}
\usepackage{bm}
\usepackage{booktabs}
\usepackage{xr-hyper}
\usepackage[colorlinks=true,allcolors=blue]{hyperref}
\usepackage{cite}

% Equation, section, table, and proposition numbers of the main text are
% read from main.aux and cited with the prefix M-, e.g. \eqref{M-eq:closure}.
\externaldocument[M-]{main}

\renewcommand{\thesection}{S\arabic{section}}
\renewcommand{\theequation}{S\arabic{equation}}
\renewcommand{\thetable}{S\Roman{table}}
\renewcommand{\thefigure}{S\arabic{figure}}

\theoremstyle{plain}
\newtheorem{proposition}{Proposition}
\renewcommand{\theproposition}{S\arabic{proposition}}

\begin{document}

\title{Supplementary Material for ``Constructing Machine Surrogates for Power System Transient Prediction from Limited Measurements: A Physics-Prior Neural ODE Approach''}

\author{Ruizhi~Yu, Wei~Gu, Shuai Lu, Ke Wang, Suhan Zhang}

\maketitle

This document supplements the main text.
Section~\ref{sec:family} lists the synchronous machine models of the family in full.
Section~\ref{sec:proof} proves Proposition~\ref{M-prop:bounded} of the main text.
Section~\ref{sec:order} shows that the learned rotor model represents the reduced-order members of the family.
Section~\ref{sec:grad} gives the gradient of the port-only loss.
Section~\ref{sec:ndf} gives the numerical differentiation formula and its residual equation.
Section~\ref{sec:params} lists the parameters of the IEEE 14-bus case study.
Equation, section, table, and proposition numbers without the prefix S refer to the main text.

\section{The Synchronous Machine Model Family}\label{sec:family}

\subsection{The Sauer--Pai Model with GENROU Saturation}\label{ssec:sp}

The sixth-order Sauer--Pai model \cite{sauer2018} with the GENROU saturation \cite{cui2021} reads
\begin{subequations}\label{eq:sm}
\begin{align}
&V_d=V_h\sin(\delta-\theta_h),\qquad V_q=V_h\cos(\delta-\theta_h),\label{eq:park}\\
&\dot\delta=\Omega_b(\omega-1),\qquad
2H\dot\omega=T_m-T_e-D(\omega-1),\label{eq:swing}\\
&0=V_d+R_aI_d-E_d''-X_q''I_q,\label{eq:stator}\\
&0=V_q+R_aI_q-E_q''+X_d''I_d,\label{eq:stator-q}\\
&T_e=E_d''I_d+E_q''I_q+(X_q''-X_d'')I_dI_q,\label{eq:torque}\\
&T_{d0}'\dot E_q'=-E_q'-(X_d-X_d')\bigl(I_d-\gamma_d\eta_d\bigr)\notag\\
&\qquad\qquad\quad+E_{fd}-S_eE_q'',\label{eq:sp}\\
&T_{q0}'\dot E_d'=-E_d'+(X_q-X_q')\bigl(I_q-\gamma_q\eta_q\bigr)-S_eg_{qd}E_d'',\label{eq:sp-ed}\\
&T_{d0}''\dot\psi_{1d}=-\psi_{1d}+E_q'-(X_d'-X_l)I_d,\label{eq:sp-p1d}\\
&T_{q0}''\dot\psi_{2q}=-\psi_{2q}-E_d'-(X_q'-X_l)I_q,\label{eq:sp-p2q}\\
&E_q''=a_dE_q'+b_d\psi_{1d},\qquad
E_d''=a_qE_d'-b_q\psi_{2q},\label{eq:sp-emf}\\
&S_e=\frac{B_s(\psi_a-A_s)^2}{\psi_a},\qquad
\psi_a=\sqrt{(E_d'')^2+(E_q'')^2}.\label{eq:sat}
\end{align}
\end{subequations}
Equations \eqref{eq:park}--\eqref{eq:torque} and \eqref{eq:sp-emf} are the swing and stator equations \eqref{M-eq:sm} of the main text, and the symbols there keep their meaning.
The rotor equations \eqref{eq:sp}--\eqref{eq:sp-p2q} take the transient EMFs $E_q',E_d'$ and the damper flux linkages $\psi_{1d},\psi_{2q}$ as rotor states, where $E_{fd}$ is the field voltage, $X_d,X_q$ the synchronous reactances, $X_d',X_q'$ the transient reactances, $X_l$ the leakage reactance, and $T_{d0}',T_{q0}',T_{d0}'',T_{q0}''$ the open-circuit time constants.
The coupling terms carry the sub-transient leakage fluxes
\begin{equation}\label{eq:eta}
\eta_d=\psi_{1d}+(X_d'-X_l)I_d-E_q',\quad
\eta_q=\psi_{2q}+(X_q'-X_l)I_q+E_d',
\end{equation}
and the constant coefficients
\begin{equation}\label{eq:gamma}
\gamma_d=\frac{X_d'-X_d''}{(X_d'-X_l)^2},\qquad
\gamma_q=\frac{X_q'-X_q''}{(X_q'-X_l)^2}.
\end{equation}
The EMF map \eqref{eq:sp-emf} has the coefficients
\begin{equation}\label{eq:ab}
\begin{aligned}
a_d&=\frac{X_d''-X_l}{X_d'-X_l},\qquad
b_d=\frac{X_d'-X_d''}{X_d'-X_l},\\
a_q&=\frac{X_q''-X_l}{X_q'-X_l},\qquad
b_q=\frac{X_q'-X_q''}{X_q'-X_l}.
\end{aligned}
\end{equation}
The saturation factor \eqref{eq:sat} is a quadratic function of the air-gap flux magnitude $\psi_a$ that vanishes for $\psi_a\le A_s$, and it enters the $q$ axis through $g_{qd}=(X_q-X_l)/(X_d-X_l)$.
Its coefficients $A_s$ and $B_s$ are fitted so that $S_e$ passes through the standard saturation points $(1.0,S_{1.0})$ and $(1.2,S_{1.2})$ \cite{cui2021}, and $S_e=0$ recovers the unsaturated Sauer--Pai model.

\subsection{The Marconato and the Anderson--Fouad Models}\label{ssec:marconato}

The Marconato model \cite{milano2010} takes the sub-transient EMFs themselves as rotor states, $E_q',E_d',E_q'',E_d''$, and its rotor equations read
\begin{subequations}\label{eq:marconato}
\begin{align}
&T_{d0}'\dot E_q'=-E_q'-(X_d-X_d'-\gamma_d^{\mathrm{M}})I_d+\Bigl(1-\frac{T_{AA}}{T_{d0}'}\Bigr)E_{fd},\\
&T_{q0}'\dot E_d'=-E_d'+(X_q-X_q'-\gamma_q^{\mathrm{M}})I_q,\\
&T_{d0}''\dot E_q''=-E_q''+E_q'-(X_d'-X_d''+\gamma_d^{\mathrm{M}})I_d+\frac{T_{AA}}{T_{d0}'}E_{fd},\\
&T_{q0}''\dot E_d''=-E_d''+E_d'+(X_q'-X_q''+\gamma_q^{\mathrm{M}})I_q,
\end{align}
\end{subequations}
where $T_{AA}$ is the additional leakage time constant of the $d$ axis and
\begin{equation}\label{eq:gamma-marconato}
\gamma_d^{\mathrm{M}}=\frac{T_{d0}''X_d''}{T_{d0}'X_d'}(X_d-X_d'),\qquad
\gamma_q^{\mathrm{M}}=\frac{T_{q0}''X_q''}{T_{q0}'X_q'}(X_q-X_q').
\end{equation}
The sub-transient EMFs are states, so the EMF map \eqref{eq:sp-emf} is the identity, and the stator and torque equations \eqref{eq:stator}--\eqref{eq:torque} apply unchanged.
The Anderson--Fouad model \cite{milano2010} is \eqref{eq:marconato} with $\gamma_d^{\mathrm{M}}=\gamma_q^{\mathrm{M}}=0$ and $T_{AA}=0$, that is, four cascaded first-order equations without coupling.

\subsection{The Reduced-Order Models}\label{ssec:reduced}

The reduced-order members drop the rotor states whose time constants are separated from the electromechanical time scale \cite{milano2010,sauer2018}.
The two-axis model drops the damper windings, whose time constants $T_{d0}''$ and $T_{q0}''$ are of the order of tens of milliseconds, and keeps
\begin{subequations}\label{eq:twoaxis}
\begin{align}
&T_{d0}'\dot E_q'=-E_q'-(X_d-X_d')I_d+E_{fd},\\
&T_{q0}'\dot E_d'=-E_d'+(X_q-X_q')I_q,\\
&0=V_d+R_aI_d-E_d'-X_q'I_q,\\
&0=V_q+R_aI_q-E_q'+X_d'I_d,\\
&T_e=E_d'I_d+E_q'I_q+(X_q'-X_d')I_dI_q,
\end{align}
\end{subequations}
so the transient EMFs act behind the transient reactances.
The one-axis model further drops $E_d'$ and keeps
\begin{subequations}\label{eq:oneaxis}
\begin{align}
&T_{d0}'\dot E_q'=-E_q'-(X_d-X_d')I_d+E_{fd},\\
&0=V_d+R_aI_d-X_qI_q,\\
&0=V_q+R_aI_q-E_q'+X_d'I_d,\\
&T_e=E_q'I_q+(X_q-X_d')I_dI_q.
\end{align}
\end{subequations}
The classical model holds a constant EMF $E'$ behind the transient reactance $X_d'$, since $T_{d0}'$ is of the order of seconds and the field flux barely changes within a first-swing study, and keeps
\begin{subequations}\label{eq:classical}
\begin{align}
&0=V_d+R_aI_d-X_d'I_q,\\
&0=V_q+R_aI_q-E'+X_d'I_d,\\
&T_e=E'I_q.
\end{align}
\end{subequations}
In every member the Park transformation \eqref{eq:park} and the swing equations \eqref{eq:swing} are unchanged.

\subsection{Saturation Conventions}\label{ssec:satconv}

Even for the saturation no consensus exists.
Reference \cite{sauer2018} distributes one open-circuit saturation function of the air-gap flux over the field and damper windings, which the GENROU form \eqref{eq:sp}--\eqref{eq:sp-ed} and \eqref{eq:sat} follows, while \cite{milano2010} admits a separate saturation function per winding, realized as a three-region piecewise curve or as a polynomial fit of the saturation characteristic.
The rotor equations thus vary across references, machine types, and model orders, which is why the main text learns the rotor model instead of presuming one member.

\section{Proof of Proposition~\ref{M-prop:bounded}}\label{sec:proof}

Proposition~\ref{M-prop:bounded} of the main text bounds every trajectory of the learned rotor model \eqref{M-eq:closure}, for any network weights, by \eqref{M-eq:ultbound}.

\begin{proof}
Write \eqref{M-eq:closure} as $\dot{\bm{x}}_r=\bm{\Lambda}\bm{x}_r+\bm{w}(t)$ with $\bm{\Lambda}=\bm{T}^{-1}A$ and $\bm{w}=\bm{T}^{-1}\bigl(B\bm{u}_r+\bm{r}_{\bm{\theta}}\bigr)$.
The output bound of \eqref{M-eq:rtheta} gives $\lvert r_{\bm{\theta},i}\rvert\le c_i$, so $\lVert\bm{w}(t)\rVert_2\le\bar{u}+\lVert\bm{T}^{-1}\bm{c}\rVert_2$ under the input bound $\lVert\bm{T}^{-1}B\bm{u}_r(t)\rVert_2\le\bar{u}$.
Every eigenvalue of $\bm{\Lambda}$ is negative, so constants $k\ge1$ and $\lambda>0$ with $\lVert e^{\bm{\Lambda}t}\rVert_2\le ke^{-\lambda t}$ exist \cite{khalil2002}.
The variation-of-constants formula gives $\bm{x}_r(t)=e^{\bm{\Lambda}t}\bm{x}_r(0)+\int_0^{t}e^{\bm{\Lambda}(t-s)}\bm{w}(s)\,\mathrm{d}s$.
Bounding the integrand by $ke^{-\lambda(t-s)}\lVert\bm{w}\rVert_2$ and integrating yields \eqref{M-eq:ultbound}.
\end{proof}

\section{Representation of the Reduced-Order Members}\label{sec:order}

The learned rotor model \eqref{M-eq:closure} is built on the rotor states of the sixth-order Sauer--Pai model.
This section shows that the reduced-order members of Table~\ref{M-tab:variants} lie inside the family that the surrogate realizes, so fitting their records incurs no structural error.

\begin{proposition}\label{prop:order}
Let $\mathcal{S}$ collect the machine models that the hardcoded equations \eqref{M-eq:sm} together with \eqref{M-eq:closure} and \eqref{M-eq:rtheta} of the main text realize as the learnable physical parameters take positive values and the network weights take real values.
Every reduced-order member of Table~\ref{M-tab:variants} is then either a member of $\mathcal{S}$ or approximated by members of $\mathcal{S}$ to any prescribed accuracy, as follows.
The two-axis and the one-axis models belong to $\mathcal{S}$.
For the classical model, any compact operating region $\mathcal{K}$, and any $\varepsilon>0$, there is a member of $\mathcal{S}$ whose state equations differ from those of the classical model by less than $\varepsilon$ on $\mathcal{K}$.
Started from the same state, the two models then stay within $\varepsilon\bigl(e^{\kappa t}-1\bigr)/\kappa$ of each other at time $t$ while inside $\mathcal{K}$, where $\kappa$ is a Lipschitz constant of the classical state equations on $\mathcal{K}$.
\end{proposition}

\begin{proof}
Consider first the two-axis model.
Choose $X_d''=X_d'$ and $X_q''=X_q'$, and set the network output to zero, $\bm{r}_{\bm{\theta}}=\bm{0}$.
The EMF-map coefficients \eqref{eq:ab} become $a_d=a_q=1$ and $b_d=b_q=0$, so the EMF map reduces to $E_q''=E_q'$ and $E_d''=E_d'$, and the damper states $\psi_{1d}$ and $\psi_{2q}$ no longer appear in it.
The stator and torque equations contain the rotor state only through $E_q''$ and $E_d''$, the first two rows of \eqref{M-eq:AB} contain only $E_q'$ and $E_d'$, and the one remaining path, the network input, is closed because $\bm{r}_{\bm{\theta}}=\bm{0}$.
The damper states therefore appear nowhere except in their own rows,
\begin{equation}\label{eq:damper-rows}
\begin{aligned}
T_{d0}''\dot\psi_{1d}&=-\psi_{1d}+E_q'-(X_d'-X_l)I_d,\\
T_{q0}''\dot\psi_{2q}&=-\psi_{2q}-E_d'-(X_q'-X_l)I_q.
\end{aligned}
\end{equation}
Each row of \eqref{eq:damper-rows} is stable and influences no other equation, so the two damper states are inert and removing them changes no other variable.
Substituting $E_q''=E_q'$, $E_d''=E_d'$, $X_d''=X_d'$, and $X_q''=X_q'$ into the stator and torque equations of \eqref{M-eq:sm} yields the interface of \eqref{eq:twoaxis}, and the remaining dynamics are the first two rows of \eqref{M-eq:AB}, namely $T_{d0}'\dot E_q'=-E_q'-(X_d-X_d')I_d+E_{fd}$ and $T_{q0}'\dot E_d'=-E_d'+(X_q-X_q')I_q$.
These are exactly the states, rotor equations, and interface of the two-axis model \eqref{eq:twoaxis}, so the two-axis model belongs to $\mathcal{S}$.

Consider next the one-axis model.
Keep the choices above and set $X_q'=X_q$ in addition.
The $E_d'$ row then loses its input term and reduces to $T_{q0}'\dot E_d'=-E_d'$.
Every simulation starts from the rotor equilibrium of Section~\ref{M-ssec:delta0} of the main text, where $E_d'=0$, and the reduced row keeps $E_d'$ at zero for all time.
With $E_d'\equiv0$ and $X_q'=X_q$, the interface of \eqref{eq:twoaxis} becomes that of \eqref{eq:oneaxis}, and the only remaining rotor equation is the $E_q'$ row.
This is exactly the one-axis model, so it belongs to $\mathcal{S}$.

Consider finally the classical model.
Choose $X_d''=X_q''=X_q'=X_d'$, which makes the interface isotropic with the single reactance $X_d'$.
The classical model further requires $E_q'$ to stay at a constant value and $E_d'$ to stay at zero, which the linear rows alone do not deliver, and the network supplies it.
The correction that achieves this exactly is $r_1=E_q'+(X_d-X_d')I_d-E_{fd}$ and $r_2=-(X_q-X_d')I_q$, since substituting $r_1$ into the $E_q'$ row gives $\dot E_q'=0$, and substituting $r_2$ into the $E_d'$ row gives $T_{q0}'\dot E_d'=-E_d'$, which holds $E_d'$ at its equilibrium value of zero as before.
The classical machine admits no field-voltage input, so $E_{fd}$ is a constant and $r_1$ is a function of $E_q'$ and $I_d$ alone.
The pair $r_1,r_2$ is therefore a continuous function of the network inputs, bounded on $\mathcal{K}$, and a network of the form \eqref{M-eq:rtheta} whose bound $\bm{c}$ exceeds that of the correction approximates it uniformly on $\mathcal{K}$ \cite{hornik1991}.
To compare the two models as differential equations, solve the two stator equations for $I_d$ and $I_q$ and substitute the solution into the remaining equations.
The solution is unique and exact because the stator equations are linear in $I_d$ and $I_q$ with the positive determinant $R_a^2+X_d''X_q''$.
The two sets of state equations then differ only through the network, by the approximation error divided by the rotor time constants, so for any $\varepsilon>0$ weights exist that keep the difference below $\varepsilon$ on $\mathcal{K}$.
Both sets are smooth on the compact $\mathcal{K}$, hence Lipschitz, and the comparison bound for two differential equations whose right-hand sides differ by at most $\varepsilon$ gives the stated deviation from a common initial state \cite{khalil2002}.
The terminal current is a smooth function of the state and the terminal voltage, so it deviates proportionally.
\end{proof}

The proposition concerns what the surrogate can represent, not what training will find.
The port records of a reduced-order machine can be fitted with zero structural error for the two-axis and one-axis members, and with arbitrarily small trajectory error for the classical member.
The parameter values in the proof are one such representation, not the only one.
The damper time constants are arbitrary once the damper states decouple, and the correction $\bm{r}_{\bm{\theta}}$ can compensate a departure of the linear part, so many combinations of parameters and weights produce the same terminal behavior.
Training therefore recovers the dynamics of a reduced-order machine without recovering the particular parameter values of the proof, and Section~\ref{M-ssec:exp_orders} of the main text confirms the recovery across the order spectrum at the level of the port behavior.
By the same degeneracy, $X_d''$ and $X_q''$ are not identifiable on a reduced-order machine, and the surrogate claims no parameter recovery for them.

\section{Gradient of the Port-Only Loss}\label{sec:grad}

The gradient of the loss \eqref{M-eq:loss} is obtained by first discretizing the forward simulation and then differentiating the discrete solution, which is exactly what reverse-mode automatic differentiation computes.
Let $\bm{y}=[\delta,\omega,\bm{x}_r^\top]^\top$ collect the surrogate state and let $\bm{\Phi}_n$ denote one step of a fixed-step fourth-order Runge--Kutta scheme applied to the surrogate dynamics under the recorded inputs.
The forward propagation of one record of $L$ steps is
\begin{equation}\label{eq:forward}
\bm{y}_0=\bm{\iota}(\bm{X},\bm{P}),\quad
\bm{y}_{n+1}=\bm{\Phi}_n(\bm{y}_n;\bm{P}),\quad n=0,\dots,L-1,
\end{equation}
where $\bm{\iota}$ is the initial-condition map of Section~\ref{M-ssec:delta0} of the main text, $\bm{X}$ are the latent initial angles, and $\bm{P}$ are the shared parameters.
For the backward propagation, write $\ell_n$ for the contribution of sample $n$ to \eqref{M-eq:loss}, and $\bm{a}_n=\partial\mathcal{L}/\partial\bm{y}_n$ for the sensitivity of the loss to the state at sample $n$.
The sensitivity obeys the reverse recursion
\begin{equation}\label{eq:adjoint}
\bm{a}_L=\frac{\partial\ell_L}{\partial\bm{y}_L},\qquad
\bm{a}_n=\frac{\partial\ell_n}{\partial\bm{y}_n}
+\Bigl(\frac{\partial\bm{\Phi}_n}{\partial\bm{y}_n}\Bigr)^{\!\top}\bm{a}_{n+1}
\end{equation}
for $n=L-1,\dots,0$.
The gradient with respect to the learnable variables $\bm{\vartheta}=(\bm{X},\bm{P})$ then assembles from three terms,
\begin{equation}\label{eq:grad}
\frac{\mathrm{d}\mathcal{L}}{\mathrm{d}\bm{\vartheta}}
=\sum_{n=0}^{L-1}\Bigl(\frac{\partial\bm{\Phi}_n}{\partial\bm{\vartheta}}\Bigr)^{\!\top}\bm{a}_{n+1}
+\sum_{n=0}^{L}\frac{\partial\ell_n}{\partial\bm{\vartheta}}
+\Bigl(\frac{\partial\bm{y}_0}{\partial\bm{\vartheta}}\Bigr)^{\!\top}\bm{a}_0,
\end{equation}
summed over the records.
The first term carries the influence of $\bm{P}$ on the dynamics, and the second its direct influence on the predicted current through the stator equations \eqref{M-eq:stator}--\eqref{M-eq:stator-q}.
The third carries the influence of the initial condition.
It is the only term that reaches $\bm{X}$, and it adds to the gradient that the first two terms already give $\bm{P}$.

\section{The Numerical Differentiation Formula}\label{sec:ndf}

Write $\bm{z}=[\bm{x}^\top,\bm{y}^\top]^\top$ for the full unknown vector of the DAE \eqref{M-eq:dae}, $\bm{F}=[\bm{f}^\top,\bm{g}^\top]^\top$ for its right-hand side, and $M=\operatorname{diag}(I,0)$ for the mass matrix, so that the DAE is $M\dot{\bm{z}}=\bm{F}(t,\bm{z})$.
The step from $t_n$ to $t_{n+1}$ is taken by the numerical differentiation formula (NDF) of order $k$ \cite{klopfensteinNumericalDifferentiationFormulas1971},
\begin{equation}\label{eq:ndf}
M\Bigl(\sum_{m=1}^{k}\frac{1}{m}\nabla^m\bm{z}_{n+1}-\kappa\gamma_k\nabla^{k+1}\bm{z}_{n+1}\Bigr)
=h\bm{F}(t_{n+1},\bm{z}_{n+1}),
\end{equation}
where $h$ is the step size, $\nabla$ is the backward difference operator with $\nabla^0\bm{z}_n=\bm{z}_n$ and $\nabla^{j+1}\bm{z}_n=\nabla^j\bm{z}_n-\nabla^j\bm{z}_{n-1}$, $\gamma_k=\sum_{j=1}^{k}1/j$, and $\kappa$ is the damping coefficient that lowers the leading truncation term.
With the predictor $\bm{z}_{n+1}^{(0)}=\sum_{m=0}^{k}\nabla^m\bm{z}_n$, equation \eqref{eq:ndf} becomes the residual equation \eqref{M-eq:ndf-resid} of the main text,
\begin{equation}\label{eq:ndf-resid}
\begin{aligned}
\bm{G}(\bm{z})&=M\bigl(\bm{z}-\bm{z}_{n+1}^{(0)}+\bm{\Psi}\bigr)-\alpha_k\bm{F}(t_{n+1},\bm{z})=\bm{0},\\
\alpha_k&=\frac{h}{(1-\kappa)\gamma_k},\qquad
\bm{\Psi}=\frac{1}{(1-\kappa)\gamma_k}\sum_{m=1}^{k}\gamma_m\nabla^m\bm{z}_n,
\end{aligned}
\end{equation}
which is one nonlinear system per step.
The simplified Newton iteration that solves it is \eqref{M-eq:simp-newton} of the main text.

\section{Parameters of the IEEE 14-Bus Case Study}\label{sec:params}

The five members of the family in Section~\ref{M-ssec:exp_orders} of the main text share the parameters of the machine at bus 2 of the IEEE 14-bus system \cite{milano2010}, namely $H=3.924$~s, $D=1.2$, $R_a=0.0052$, $X_d=1.75$, $X_q=1.633$, $X_d'=0.308$, $X_q'=0.6$, $X_d''=X_q''=0.217$, and $X_l=0$ in per unit, and $T_{d0}'=6.1$, $T_{q0}'=0.3$, $T_{d0}''=0.04$, and $T_{q0}''=0.099$~s.
The GENROU setting adds the saturation factors $S_{1.0}=0.27$ and $S_{1.2}=1.14$.
Every learnable parameter of the surrogate starts from these exact values.

\bibliographystyle{IEEEtran}
\bibliography{refs}

\end{document}
